\chapter{Composant Machine Learning}

\section{Dataset et Déséquilibre}

Le projet utilise le dataset Kaggle Credit Card Fraud Detection, composé d'environ 284\,807 transactions dont 492 fraudes. La classe positive représente environ 0,17\% des observations, ce qui impose une stratégie adaptée aux classes fortement déséquilibrées.

Les variables \repoPath{V1} à \repoPath{V28} sont anonymisées par PCA. Les variables lisibles sont \repoPath{Time}, \repoPath{Amount} et \repoPath{Class}.

\section{Préparation des Données}

Le pipeline suit les étapes suivantes:

\begin{enumerate}
    \item chargement du dataset;
    \item nettoyage et suppression des doublons;
    \item clipping des montants extrêmes selon la logique de preprocessing;
    \item feature engineering;
    \item split stratifié train, validation et test;
    \item scaling et SMOTE uniquement dans les folds de cross-validation via \repoPath{imblearn.Pipeline}.
\end{enumerate}

Le point méthodologique essentiel est l'absence de data leakage: le scaler et SMOTE ne sont pas ajustés avant le split, et le test set reste intouché jusqu'à l'évaluation finale.

\section{Feature Engineering}

Le projet ajoute 9 variables dérivées aux 30 variables initiales, pour un total de 39 features attendues par le modèle. Les principales familles sont:

\begin{itemize}
    \item variables temporelles: heure transactionnelle et période de journée;
    \item variables de montant: log, bins et interactions;
    \item statistiques sur les composantes PCA: moyenne et écart-type;
    \item score de risque composite construit à partir de variables corrélées à la fraude.
\end{itemize}

\section{Modèles Comparés}

Le script \repoPath{src/core/models/train.py} compare quatre familles:

\begin{itemize}
    \item Logistic Regression;
    \item Decision Tree;
    \item Random Forest;
    \item XGBoost.
\end{itemize}

La sélection se fait sur validation, en tenant compte de la PR-AUC, du F1-score validation et du F1-score de cross-validation. Le modèle final est ré-entraîné sur train + validation, puis évalué une seule fois sur le test set.

\section{Résultats}

\begin{table}[htbp]
\centering
\caption{Matrice de confusion finale}
\begin{tabular}{lrr}
\toprule
 & Prédit normal & Prédit fraude\\
\midrule
Réel normal & 42\,485 & 3\\
Réel fraude & 15 & 56\\
\bottomrule
\end{tabular}
\end{table}

\begin{figure}[htbp]
\centering
\includegraphics[width=0.82\textwidth]{../fraud-detection/figures/feature_importance.png}
\caption{Importance des variables issue des analyses du projet}
\end{figure}

\begin{figure}[htbp]
\centering
\includegraphics[width=0.82\textwidth]{../fraud-detection/figures/threshold_optimization.png}
\caption{Analyse du seuil de décision}
\end{figure}

\section{Artefacts}

Les artefacts principaux sont:

\begin{itemize}
    \item \repoPath{models/best_model.joblib}: pipeline complet sauvegardé;
    \item \repoPath{models/scaler.joblib}: scaler de secours;
    \item \repoPath{models/feature_columns.json}: colonnes attendues en prédiction;
    \item \repoPath{models/metrics.json}: métriques finales;
    \item \repoPath{models/reference_stats.json}: statistiques de référence pour monitoring si générées.
\end{itemize}
